\section{Related work}
Whole-brain sensorimotor modeling is the foundation of this system. \citet{shiu2024} use a connectome-derived spiking model to investigate feeding and grooming circuits and test selected predictions experimentally. Our assembly inherits this modeling direction and adds interfaces and persistent state. We do not treat the inherited architecture or previously established sensorimotor phenomena as new discoveries.

Connectome constraints can also be combined with task optimization. The flyvis study fits visual-system dynamics to a motion task and compares neural responses with physiological measurements \citep{lappalainen2024}. FlyGM uses a connectome-structured controller in embodied locomotion tasks \citep{jin2026}. These works address visual prediction and physical control, respectively. The current platform does not include their trained visual model or their embodied controller, and its motor-neuron outputs are not comparable to measured locomotion performance.

Neural inputs to language generation also predate this work. BrainLLM conditions generation on representations decoded from human brain recordings \citep{ye2025}. Our targets are simulator-derived captions about interventions and model outputs rather than language experienced by a human participant. The continuous-prefix interface is an implementation choice within a larger research tool, not a new neural-decoding algorithm.

Interactive repositories such as Fly Escape and DOOMFLY already connect fly-connectome simulations to game environments \citep{flyescape,doomfly}. They make a broad claim of ``the first interactive fly connectome'' inappropriate. The narrower question for this resource is whether its chemical-input, persistent-memory, motor-output, and language interfaces support useful reproducible experiments together. The recorded chemical, memory, and language cases define the scope of our resource contribution; we make no priority claim for the combination.
